\section{Patching centres and the ideal-square obstruction}\label{sec:global}

\subsection{Only finitely many local tests are needed}
First check $a\in R^\times$ and that every nonzero valuation of $b$
is divisible by $d-1$. These questions use finite factorizations of
principal fractional ideals. Let $S$ be the finite set of places
where $v(b)\ne0$ or $v(q_j)<0$ for at least one $j$. Outside $S$,
the choice $s=1,r=0$ passes. In particular, a place dividing $d$ need
not be searched merely because its residue characteristic is wild.

For $v\in S$, a scale $s_v\in K^\times$ with valuation $k_v$ can
be chosen: an element of $\mathfrak p_v\setminus\mathfrak p_v^2$
has valuation one, and an integer power gives the required scale.
This does not require the prime ideal to be principal. Representatives
of $\OO_{K_v}/d\OO_{K_v}$ can be taken in $R$, using
\[
 R/\mathfrak p_v^{e_v}\simeq
 \OO_{K_v}/\mathfrak p_v^{e_v}\OO_{K_v}
 =\OO_{K_v}/d\OO_{K_v}.
\]
Consequently all tested centres $r_v=r_0+(s_v/d)\beta$ can be taken
in $K$. If every place passes, set $r_v=0,s_v=1$ outside $S$.
Then $I$ in \eqref{eq:scaleideal} has finite support and its prescribed
valuations immediately give $(b)I^{d-1}=R$.

\begin{lemma}[Additive patching]\label{lem:patch}
For these local discs $r_v+I\OO_{K_v}$ there exists $r\in K$ such
that $r-r_v\in I\OO_{K_v}$ at every finite place. The set of all
such $r$ is one coset of $I$ in $K$.
\end{lemma}
\begin{proof}
Choose $0\ne D\in R$ clearing denominators so that $J=DI$ is an
integral ideal and every $Dr_v$ is integral at every place. This is
possible since only finitely many $r_v$ are nonzero elements of $K$.
Write $n_v=v(J)\ge0$. For the finitely many places with $n_v>0$,
the local residue $Dr_v$ modulo $\mathfrak p_v^{n_v}$ has a
representative in $R$. The Chinese remainder theorem supplies
$z\in R$ satisfying
\[
 z\equiv Dr_v\pmod{\mathfrak p_v^{n_v}\OO_{K_v}}
 \qquad(n_v>0).
\]
This includes any new places introduced by clearing denominators;
it is not restricted to $S$. When $n_v=0$, both terms are already
integral. Thus $z-Dr_v\in J\OO_{K_v}$ at every place, and $r=z/D$
has the required property. Two solutions differ by an element of
$K\cap\bigcap_v I\OO_{K_v}=I$, by the valuation description of
a fractional ideal. Conversely every element of that coset works.
\end{proof}

\subsection{Necessity of a principal square}
Assume now that $T(z)=Az+c$ is globally good and choose $r$ as in
Lemma~\ref{lem:patch}. Local uniqueness gives
\begin{equation}\label{eq:localmodule}
 A\OO_{K_v}^2=I\OO_{K_v}\oplus I\OO_{K_v},\qquad
 c-(r,r)\in I\OO_{K_v}\oplus I\OO_{K_v}.
\end{equation}
These imply $AR^2=I\oplus I$ and $c-(r,r)\in I\oplus I$.
For the first equality, one inclusion follows by localization.
If $x\in I\oplus I$, then $A^{-1}x\in\OO_{K_v}^2$ for every
finite $v$, hence $A^{-1}x\in R^2$, giving the other inclusion.
Taking second exterior powers in $\bigwedge^2K^2\simeq K$ yields
\begin{equation}\label{eq:detideal}
 I^2=(\det A).
\end{equation}
This is the necessary determinant obstruction. In particular it
does not assert that $I$ itself must be principal.

\subsection{An explicit global basis}
We include the needed Dedekind-module construction to keep the
sufficiency argument valid for nonprincipal ideals.

\begin{lemma}[Two generators]\label{lem:twogen}
Every nonzero fractional ideal of $R$ has two $R$-generators.
\end{lemma}
\begin{proof}
Clear denominators and assume $I$ is integral. Choose
$0\ne\alpha\in I$ and let $P$ be the finite set of prime ideals
where $v(\alpha)>v(I)$. For each $\mathfrak p\in P$ choose
$\beta_{\mathfrak p}\in I\setminus\mathfrak p I$.
Such an element exists since the invertible ideal $I$ has
one-dimensional nonzero quotient $I/\mathfrak p I$ over
$R/\mathfrak p$. The module version of the Chinese remainder
theorem gives $\beta\in I$ congruent to
$\beta_{\mathfrak p}$ modulo $\mathfrak p I$ for all
$\mathfrak p\in P$. At such primes $v(\beta)=v(I)$; at all
other primes $v(\alpha)=v(I)$. Thus $(\alpha,\beta)$ and $I$
have the same valuations and are equal. If $P$ is empty, take
$\beta=0$. Restoring denominators proves the fractional case.
\end{proof}

\begin{lemma}[Basis from a principal square]\label{lem:basis}
If $I^2=(\gamma)$, there is an explicit matrix $A\in\GL_2(K)$
with $AR^2=I\oplus I$.
\end{lemma}
\begin{proof}
Write $I=(\alpha,\beta)$ by Lemma~\ref{lem:twogen}. From
$II^{-1}=R$ choose $u,v\in I^{-1}$ with $\alpha u+\beta v=1$.
Set
\begin{equation}\label{eq:basis}
 A=\begin{pmatrix}\alpha&-\gamma v\\
                    \beta&\gamma u\end{pmatrix}.
\end{equation}
All entries lie in $I$ because $\gamma I^{-1}=I$, and
$\det A=\gamma$. For $x,y\in I$,
\begin{equation}\label{eq:basisinverse}
 A^{-1}\binom xy
 =\binom{ux+vy}{(-\beta x+\alpha y)/\gamma}\in R^2.
\end{equation}
The first entry is integral by $I^{-1}I=R$, and the second by
$I^2=(\gamma)$. This proves both inclusions in $AR^2=I\oplus I$.
\end{proof}

The two lemmas and \eqref{eq:detideal} prove
\eqref{eq:obstruction}, including the equivalence with freeness.
For sufficiency in the dynamical theorem, take the matrix
\eqref{eq:basis} and $T(z)=Az+(r,r)$. At a finite place, use the
previous passing pair $(s_v,r_v)$. Then
\begin{equation}\label{eq:localfactor}
 S_{s_v,r_v}^{-1}T(z)
 =\frac{A}{s_v}z+\frac{r-r_v}{s_v}(1,1).
\end{equation}
Here $A/s_v\in\GL_2(\OO_{K_v})$ by equality of the local lattices,
and the translation is integral by Lemma~\ref{lem:patch}.
Lemma~\ref{lem:rigidity} therefore makes $T$ good at every place.
The coefficients of $T^{-1}FT$ and its inverse lie in
$K\cap\bigcap_v\OO_{K_v}=R$.

\subsection{All coordinate maps and choice independence}
The necessity argument proved that every global map satisfies
\eqref{eq:allglobal}. Conversely, for every $A,c$ satisfying that
formula the same local factorization \eqref{eq:localfactor}, with
$c-(r_v,r_v)$ in its translation part, is integral and invertible.
Hence it gives a good model. If $T_1,T_2$ satisfy the formula,
$T_1^{-1}T_2$ has linear part in $\GL_2(R)$ and translation in
$R^2$, because it carries $R^2$ bijectively to itself. This proves
the global left-coset statement.

Changing $s_v$ by a unit preserves the scale ideal; changing a passing
centre within $s_v\OO_{K_v}$ preserves its disc; changing the patched
$r$ within $I$ preserves \eqref{eq:allglobal}. The ideal $I$ itself
is uniquely prescribed by the valuations of $b$. If the original
map is expressed in another global affine coordinate system, its
unique local lattice is carried by that coordinate change; its
determinant ideal changes by a principal factor. Thus the vanishing
of the obstruction is also invariant under such a reexpression.
This completes the proof of Theorem~\ref{thm:main}.

\paragraph{A terminating arithmetic prescription.}
For clarity, the constructive content is the following finite rule.
\begin{enumerate}
\item Factor $(a)$ and $(b)$. Reject unless $a\in R^\times$ and
$(d-1)\mid v(b)$ at every finite place.
\item At each place in the finite set $S$, choose a scale of the
required valuation and enumerate \eqref{eq:candidates}. Test all
coefficients by \eqref{eq:coefftest0}--\eqref{eq:coefftest}; reject
if no candidate passes.
\item Form $I$ and patch $r$ using Lemma~\ref{lem:patch}. Test
whether $I^2$ is principal. If not, reject the global model problem.
\item If $I^2=(\gamma)$, obtain two generators, solve
$\alpha u+\beta v=1$ in $I^{-1}$ and use \eqref{eq:basis}.
Formula~\eqref{eq:allglobal} then describes every answer.
\end{enumerate}
The prescription uses exact number-field ideal arithmetic and finite
residue rings. It states no practical complexity bound and was not
implemented or executed as part of this proof.
